\appendix

\section{Test Vectors}
\label{sec:appendix-vectors}

\subsection*{Vector 1: empty content, deterministic header}

\begin{itemize}
\item \textbf{Content:} \texttt{b""} (zero bytes)
\item \textbf{Timestamp (ms):} \texttt{1700000000000}
\item \textbf{Counter:} \texttt{0}
\item \textbf{Tail (SHAKE-256(\texttt{b""}, 64 bytes)):} the canonical FIPS 202 \citep{fips202} SHAKE-256 output for the empty string, truncated to 64 bytes.
\end{itemize}

The 73-byte binary is \texttt{TS(6, big-endian) || 0x03 || 0x00 0x00 || tail(64)}. The exact hex bytes for this vector live in \texttt{tessera-hash/conformance/vectors/tessera\_hash\_hash\_only.json}.

\subsection*{Vector 2: short content}

\begin{itemize}
\item \textbf{Content:} \texttt{b"abc"}
\item \textbf{Tail:} SHAKE-256(\texttt{b"abc"}, 64 bytes), canonical FIPS 202 output.
\end{itemize}

Used by the round-trip tests in \texttt{tessera\_hash\_verify.json}: a one-bit flip anywhere in the content changes the tail in many bit positions (avalanche property of Keccak-f[1600]), so \texttt{verify(raw, modified\_content)} returns \texttt{false}.

\section{Quick Reference}
\label{sec:appendix-quickref}

\begin{center}
\small
\begin{tabular}{@{}ll@{}}
\toprule
\textbf{Property} & \textbf{Value} \\
\midrule
Wire size                 & 73 bytes \\
Header size               & 9 bytes \\
Tail size                 & 64 bytes (512 bits) \\
Tail primitive            & SHAKE-256 (FIPS 202) \\
Timestamp field           & 48-bit big-endian ms, offset 0 \\
Variant byte              & \code{0x03} at offset 6 \\
Counter field             & 16-bit big-endian, offset 7 \\
Encoding length: Base64URL & 98 chars, no padding \\
Encoding length: Crockford & 117 chars, case-insensitive on decode \\
Classical collision bound & $2^{256}$ \\
Quantum collision bound (BHT) & $\sim 2^{170}$ \\
Classical second-preimage & $2^{512}$ \\
Quantum second-preimage (Grover) & $2^{256}$ \\
Min sanity timestamp      & 1{,}700{,}000{,}000{,}000 ms (TesseraID-parity) \\
\bottomrule
\end{tabular}
\end{center}

\section{Glossary}
\label{sec:appendix-glossary}

\textbf{BHT (Brassard--H{\o}yer--Tapp).} Quantum collision-finding algorithm \citep{bht1998}. Costs $\sim 2^{n/3}$ operations for an $n$-bit space.

\textbf{Birthday bound.} Classically $\sqrt{\pi/2} \cdot 2^{n/2}$ generations before an expected collision in an $n$-bit random space.

\textbf{Content identifier.} An identifier derived deterministically from content (\tesserahash{}, IPFS CID, Git commit hash) as opposed to a random unique identifier (\tesseraid{}, UUIDv4, ULID).

\textbf{Crockford Base32.} Douglas Crockford's Base32 alphabet (\code{0--9} then \code{A--Z} minus \code{I/L/O/U}) \citep{crockford}. Case-insensitive on decode; ASCII-sort-preserving.

\textbf{Grover's algorithm.} Quantum search reducing $O(N)$ to $O(\sqrt{N})$ \citep{grover1996}.

\textbf{HMAC.} Keyed hash construction \citep{rfc2104}. \tesserahash{} is not an HMAC; see Section~\ref{sec:security}.

\textbf{KAT (Known-Answer Test).} A test vector pinning an exact expected output for a given input. \tesserahash{}'s KAT vectors live in \texttt{conformance/vectors/}.

\textbf{Keccak-f[1600].} The permutation underlying SHA-3 and SHAKE \citep{fips202}.

\textbf{SHAKE-256.} The variable-output-length member of the SHA-3 family \citep{fips202}, with a 256-bit security level. TesseraHash uses it with a fixed 64-byte (512-bit) output.

\textbf{Variant byte.} Byte at offset 6 of a \tesserahash{}. Always \code{0x03} for v1.

\section{The \texorpdfstring{\tessera{}}{Tessera} Monorepo}
\label{sec:appendix-monorepo}

For the cross-project relationship (variant-byte allocation across \tesseraid{}/\tesserahash{}/\tesseraseal{}, the canonical \tesseraid{}-as-AD and \tesserahash{}-as-AD patterns, shared \texttt{common/} library code, shared conformance harness), see Section~\ref{sec:monorepo}. The summary of cross-project byte allocations:

\begin{center}
\small
\begin{tabular}{@{}llll@{}}
\toprule
\textbf{Byte} & \textbf{Owner} & \textbf{Meaning} & \textbf{Status} \\
\midrule
\code{0x00}     & ---        & Permanently reserved (never assign) & MUST reject \\
\code{0x01}     & \tesseraid{}    & Standard generator                    & assigned \\
\code{0x02}     & \tesseraid{}    & HLC generator                         & assigned \\
\code{0x03}     & \tesserahash{}  & SHAKE-256 / 64-byte tail              & assigned (this document) \\
\code{0x04}     & \tesseraseal{} & Encryption envelope (suite byte at offset 7) & assigned \\
\code{0x05}--\code{0xFE} & --- & Reserved-future                  & MUST reject \\
\code{0xFF}     & ---        & Permanently reserved (never assign) & MUST reject \\
\bottomrule
\end{tabular}
\end{center}

\section{Author and Document Provenance}
\label{sec:appendix-provenance}

\textbf{Author:} Bahram Samani, \tessera{} monorepo maintainer.

This document is built from LaTeX sources under \texttt{tessera-hash/latex/} (\texttt{main.tex} + \texttt{secNN\_*.tex}); the build chain is \texttt{pdflatex $\rightarrow$ bibtex $\rightarrow$ pdflatex $\times$ 2} per the Makefile.

\textbf{Authority order} (when sources disagree): this whitepaper, then the Python reference at \texttt{tessera-hash/python/tessera\_hash.py}, then the canonical C ABI at \texttt{tessera-hash/c/tessera\_hash.h}, then every other language. A cross-language disagreement is never resolved by changing the spec to match an outlier; the outlier is fixed.

\textbf{Hard invariants} (from \texttt{tessera-hash/CLAUDE.md}):

\begin{itemize}
\item \textbf{Hash primitive: SHAKE-256.} Keccak-f[1600] sponge with arbitrary output length, used here with fixed 64-byte (512-bit) output.
\item \textbf{Vendored Keccak.} C/C++/Java implementations use the shared Keccak source under \texttt{common/}; no project re-vendors a second copy.
\item \textbf{Canonical encodings.} Base64URL (98 chars, no padding) and Crockford Base32 (117 chars, case-insensitive on decode, \code{I/L $\to$ 1} and \code{O $\to$ 0} confusables accepted on decode). Encoded lengths are fixed per variant.
\item \textbf{Variant discipline.} \code{0x03} only; any unknown variant byte is \code{TESSERA\_HASH\_E\_UNKNOWN\_VAR}. Cross-project bytes (\code{0x01}, \code{0x02}, \code{0x04}) and reserved-future bytes are rejected.
\item \textbf{Generator locking.} The generator is internally locked; the \texttt{(last\_ts, counter)} read-modify-write is atomic.
\end{itemize}

\textbf{Conformance gate:} \texttt{tools/conformance/run.sh} is the branch-protection check on \texttt{main}; \tesseraid{}, \tesserahash{}, and \tesseraseal{} flow through the same harness. TesseraHash's KAT count: 16 \texttt{decode}, 8 \texttt{hash\_only}, encode/decode round-trip pairs for both encodings, \texttt{verify} positive/negative cases, \texttt{invalid} negative cases. Latest cross-language run (2026-05-21) reports 0 failures across the alpha languages.
